\documentclass[aspectratio=169]{beamer}

\usepackage[utf8]{inputenc}
\usepackage[T1]{fontenc}
\usepackage[activeacute]{babel}
\usepackage{amsmath, amssymb}
\usepackage{booktabs}
\usepackage{xcolor}
\usepackage{tcolorbox}

\usetheme{Madrid}
\usecolortheme[RGB={30,60,120}]{structure}

\setbeamertemplate{navigation symbols}{}
\setbeamertemplate{footline}{
  \leavevmode
  \hbox{
    \begin{beamercolorbox}[wd=.45\paperwidth,ht=2.5ex,dp=1ex,left,leftskip=1em]{author in head/foot}
      \usebeamerfont{author in head/foot}\insertshortauthor
    \end{beamercolorbox}
    \begin{beamercolorbox}[wd=.35\paperwidth,ht=2.5ex,dp=1ex,center]{title in head/foot}
      \usebeamerfont{title in head/foot}MAD1 -- Unidad 5
    \end{beamercolorbox}
    \begin{beamercolorbox}[wd=.2\paperwidth,ht=2.5ex,dp=1ex,right,rightskip=1em]{date in head/foot}
      \usebeamerfont{date in head/foot}\insertframenumber{} / \inserttotalframenumber
    \end{beamercolorbox}
  }
}

\definecolor{unsazul}{RGB}{30,60,120}
\definecolor{unsazulclaro}{RGB}{220,230,245}
\definecolor{verdeapp}{RGB}{20,120,60}

\tcbuselibrary{skins, breakable}

\newtcolorbox{aplicacion}{
  colback=verdeapp!8,
  colframe=verdeapp,
  fonttitle=\bfseries\small,
  title={Aplicaciones reales},
  left=4pt, right=4pt, top=4pt, bottom=4pt
}

\newtcolorbox{concepto}{
  colback=unsazulclaro,
  colframe=unsazul,
  fonttitle=\bfseries\small,
  title={Idea clave},
  left=4pt, right=4pt, top=4pt, bottom=4pt
}

\newtcolorbox{atencion}{
  colback=orange!10,
  colframe=orange!80!black,
  fonttitle=\bfseries\small,
  title={Atenci\'on},
  left=4pt, right=4pt, top=4pt, bottom=4pt
}

\title[Unidad 5 -- Autoencoders]{Unidad 5: Aprendizaje No Supervisado}
\subtitle{Bloque 4 -- Aprendizaje profundo no supervisado}
\author[MAD1 -- UNS]{Métodos de Análisis de Datos I}
\institute{Departamento de Matemática -- Universidad Nacional del Sur}
\date{}

%---------------------------------------------------------------
\begin{document}
%---------------------------------------------------------------

\begin{frame}
\titlepage
\end{frame}

%---------------------------------------------------------------
\begin{frame}{¿Por qué redes neuronales no supervisadas?}
\begin{columns}[T]
\column{0.55\textwidth}
Los métodos clásicos (PCA, KDE, K-means) tienen limitaciones cuando la estructura de los datos es \textbf{no lineal} y de muy alta dimensión.

\medskip
\textbf{¿Qué aportan las redes neuronales?}
\begin{itemize}
  \item Pueden capturar relaciones \textbf{no lineales} complejas.
  \item Escalan a dimensiones muy altas (imágenes, audio, texto).
  \item Aprenden representaciones jerárquicas de los datos.
\end{itemize}

\medskip
\begin{concepto}
El principio fundamental: entrenar la red para que \textbf{comprima y reconstruya} los datos. Lo que la red aprende a preservar en la compresión es la estructura relevante. No necesitamos etiquetas: la señal de aprendizaje es el error de reconstrucción.
\end{concepto}

\column{0.42\textwidth}
\textbf{Autoencoder en este bloque:}

\medskip
{\small
\begin{tabular}{ll}
\toprule
\textbf{Componente} & \textbf{Rol} \\
\midrule
Encoder    & Comprime $\mathbf{x} \to \mathbf{z}$ \\
Bottleneck & Representación comprimida \\
Decoder    & Reconstruye $\mathbf{z} \to \hat{\mathbf{x}}$ \\
\bottomrule
\end{tabular}
}

\medskip
\begin{aplicacion}
Compresión de imágenes, detección de anomalías en manufactura, eliminación de ruido en señales médicas, representaciones para NLP, generación de nuevos datos.
\end{aplicacion}
\end{columns}
\end{frame}

%---------------------------------------------------------------
\begin{frame}{Arquitectura del autoencoder}
\begin{columns}[T]
\column{0.52\textwidth}
\[
\mathbf{x} \in \mathbb{R}^p \;\xrightarrow{\;\text{encoder}\;}\; \mathbf{z} \in \mathbb{R}^q \;\xrightarrow{\;\text{decoder}\;}\; \hat{\mathbf{x}} \in \mathbb{R}^p
\]

\medskip
\textbf{Función de pérdida:}
\[
\mathcal{L} = \frac{1}{n}\sum_{i=1}^{n}\|\mathbf{x}_i - \hat{\mathbf{x}}_i\|^2
\]

\medskip
\textbf{Entrenamiento:} backpropagation estándar. No se necesita ninguna etiqueta $y_i$.

\medskip
\textbf{Arquitectura típica para datos tabulares:}
\medskip
{\small
\begin{tabular}{ll}
\toprule
\textbf{Capa} & \textbf{Detalle} \\
\midrule
Entrada    & $p$ neuronas \\
Encoder 1  & Dense($d_1$, ReLU), $d_1 < p$ \\
Encoder 2  & Dense($q$, ReLU), bottleneck \\
Decoder 1  & Dense($d_1$, ReLU) \\
Salida     & Dense($p$, lineal o sigmoid) \\
\bottomrule
\end{tabular}
}

\column{0.45\textwidth}
\begin{concepto}
Un autoencoder es como un sistema de mensajería con canal angosto: el encoder resume la novela en un telegrama ($\mathbf{z}$) y el decoder reconstruye la novela desde el telegrama. Si el canal es angosto, la red aprende qué es lo más importante.
\end{concepto}

\medskip
\begin{aplicacion}
\textbf{Imágenes:} el bottleneck de un autoencoder entrenado en rostros aprende variables latentes como ``ángulo de la cara'', ``iluminación'', ``expresión'', aunque nadie se las indicó.

\medskip
\textbf{Texto:} autoencoders sobre embeddings de párrafos aprenden a representar el ``tema'' del texto en pocas dimensiones.
\end{aplicacion}
\end{columns}
\end{frame}

%---------------------------------------------------------------
\begin{frame}{Autoencoder como reductor de dimensionalidad}
\begin{columns}[T]
\column{0.52\textwidth}
El vector latente $\mathbf{z}$ es una representación de $\mathbf{x}$ en $q$ dimensiones.

\medskip
\textbf{Comparación con PCA:}

\medskip
{\small
\begin{tabular}{lll}
\toprule
 & \textbf{PCA} & \textbf{Autoencoder} \\
\midrule
Transformación & Lineal & No lineal \\
Computación    & Analítica (SVD) & Iterativa (GD) \\
Interpretación & Alta (loadings) & Baja \\
Escala         & Alta & Alta \\
\bottomrule
\end{tabular}
}

\medskip
\begin{concepto}
Un autoencoder \textbf{lineal} (sin activaciones no lineales) de una sola capa oculta es matemáticamente equivalente a PCA. Las activaciones no lineales son lo que agrega potencia real.
\end{concepto}

\column{0.45\textwidth}
\begin{aplicacion}
\textbf{Genómica:} reducir perfiles de expresión de 20.000 genes a 10--30 dimensiones latentes, capturando relaciones no lineales entre genes que PCA no ve.

\medskip
\textbf{Sensores industriales:} comprimir señales de vibración de 500 canales simultáneos a 5 dimensiones latentes para monitoreo en tiempo real.

\medskip
\textbf{Astronomía:} representar espectros de galaxias en pocas dimensiones para clasificación y búsqueda de objetos raros.

\medskip
\textbf{Finanzas:} comprimir miles de variables macroeconómicas a factores latentes no lineales para predicción.
\end{aplicacion}
\end{columns}
\end{frame}

%---------------------------------------------------------------
\begin{frame}{Autoencoder para detección de anomalías}
\begin{columns}[T]
\column{0.52\textwidth}
\textbf{Principio:} entrenar el autoencoder \textbf{solo con datos normales}.

\medskip
Un autoencoder bien entrenado en datos normales:
\begin{itemize}
  \item Reconstruye bien los datos normales: $\|\mathbf{x} - \hat{\mathbf{x}}\|^2$ pequeño.
  \item Reconstruye mal las anomalías: el modelo no ``aprendió'' ese patrón.
\end{itemize}

\medskip
\textbf{Score de anomalía:}
\[
\text{score}(\mathbf{x}) = \|\mathbf{x} - \hat{\mathbf{x}}\|^2
\]

Se fija un umbral $\tau$: si $\text{score}(\mathbf{x}) > \tau \Rightarrow$ anomalía.

\medskip
\begin{concepto}
El error de reconstrucción actúa como un termómetro de rareza: cuanto más difícil le resulta al autoencoder reconstruir un punto, más raro es ese punto.
\end{concepto}

\column{0.45\textwidth}
\begin{aplicacion}
\textbf{Manufactura:} entrenar con imágenes de piezas sin defectos. Las piezas con grietas o deformaciones tienen error de reconstrucción alto.

\medskip
\textbf{Fraude bancario:} entrenar con transacciones legítimas. Las transacciones fraudulentas tienen patrones no vistos que el autoencoder no puede reconstruir bien.

\medskip
\textbf{Cardiología:} entrenar con ECGs (electrocardiogramas) normales. Arritmias raras generan error de reconstrucción alto, disparando alertas automáticas.

\medskip
\textbf{Redes informáticas:} entrenar con tráfico normal. Ataques de intrusión, que tienen patrones distintos, son detectados por su alta pérdida de reconstrucción.

\medskip
\textbf{Física de partículas:} en el CERN, detectar eventos anómalos entre millones de colisiones en el LHC.
\end{aplicacion}
\end{columns}
\end{frame}

%---------------------------------------------------------------
\begin{frame}{Denoising autoencoders}
\begin{columns}[T]
\column{0.52\textwidth}
\textbf{Idea:} corromper los datos de entrada con ruido y entrenar para reconstruir la versión limpia.

\medskip
\[
\tilde{\mathbf{x}} = \mathbf{x} + \boldsymbol{\varepsilon}, \quad \boldsymbol{\varepsilon} \sim \mathcal{N}(\mathbf{0}, \sigma^2 \mathbf{I})
\]

\medskip
\textbf{Pérdida:}
\[
\mathcal{L} = \frac{1}{n}\sum_{i=1}^{n}\|\mathbf{x}_i - \text{decoder}(\text{encoder}(\tilde{\mathbf{x}}_i))\|^2
\]

\medskip
\begin{concepto}
Al forzar al encoder a ignorar el ruido y capturar solo la estructura real, el autoencoder aprende representaciones más robustas y generalizables. El ruido actúa como regularización implícita.
\end{concepto}

\column{0.45\textwidth}
\begin{aplicacion}
\textbf{Imágenes médicas:} eliminar ruido en imágenes de tomografía (CT) o resonancia magnética (MRI) adquiridas con dosis bajas de radiación.

\medskip
\textbf{Audio:} separar la voz del ruido de fondo en grabaciones con micrófono de baja calidad.

\medskip
\textbf{Astronomía:} recuperar señales débiles de galaxias lejanas contaminadas por ruido instrumental.

\medskip
\textbf{Texto:} corregir errores tipográficos o reconstruir texto con palabras faltantes en documentos digitalizados con OCR (reconocimiento óptico de caracteres) imperfecto.

\medskip
\textbf{Sensores IoT:} limpiar lecturas ruidosas de sensores de temperatura, humedad o presión en redes industriales.
\end{aplicacion}
\end{columns}
\end{frame}

%---------------------------------------------------------------
\begin{frame}{Consideraciones prácticas}
\begin{columns}[T]
\column{0.52\textwidth}
\textbf{Dimensión del bottleneck:}
\begin{itemize}
  \item Demasiado pequeña: pierde información, reconstrucción pobre.
  \item Demasiado grande: puede memorizar en lugar de generalizar.
  \item Explorar con validación cruzada o curva de reconstrucción.
\end{itemize}

\medskip
\textbf{Regularización:}
\begin{itemize}
  \item Dropout en encoder y decoder.
  \item Regularización L2 en los pesos.
  \item Denoising (ruido en la entrada).
\end{itemize}

\medskip
\textbf{Preprocesamiento:}
\begin{itemize}
  \item Escalar a $[0,1]$ si la salida es sigmoid.
  \item Estandarizar si la salida es lineal.
\end{itemize}

\medskip
\textbf{Arquitectura:} para datos tabulares, 2--3 capas ocultas son suficientes. Para imágenes, usar capas convolucionales.

\column{0.45\textwidth}
\begin{atencion}
El umbral de anomalía $\tau$ es crítico y depende del dominio. Definirlo mirando la distribución de errores de reconstrucción en el conjunto de entrenamiento (por ejemplo, percentil 95 o 99).
\end{atencion}

\medskip
\begin{aplicacion}
\textbf{CARDIO-PLACA (ejemplo propio):} entrenar un autoencoder con los 25 predictores clínicos de pacientes sanos. Los pacientes con placa de difícil diagnóstico generan mayor error de reconstrucción, pudiendo complementar al modelo de clasificación supervisado.
\end{aplicacion}

\medskip
\begin{concepto}
Los autoencoders son el punto de entrada al aprendizaje profundo no supervisado. Su extensión natural son los VAE (Autoencoders Variacionales), que generan nuevas muestras en lugar de solo reconstruir.
\end{concepto}
\end{columns}
\end{frame}

%---------------------------------------------------------------
\begin{frame}{Resumen: Bloque 4 y cierre de la Unidad 5}

\begin{columns}[T]
\column{0.48\textwidth}
\textbf{Autoencoders en síntesis:}

\medskip
{\small
\begin{tabular}{ll}
\toprule
\textbf{Uso} & \textbf{Mecanismo} \\
\midrule
Reducción de dim. & Vector latente $\mathbf{z}$ \\
Detección anomalías & Error de reconstrucción \\
Eliminación de ruido & Denoising autoencoder \\
\bottomrule
\end{tabular}
}

\medskip
\textbf{Vs.\ métodos clásicos:}
\begin{itemize}
  \item Más potente que PCA para estructuras no lineales.
  \item Más costoso computacionalmente.
  \item Menos interpretable.
  \item Requiere más datos para entrenar bien.
\end{itemize}

\column{0.48\textwidth}
\textbf{Cierre de la Unidad 5 -- los 4 bloques:}

\medskip
{\small
\begin{tabular}{ll}
\toprule
\textbf{Bloque} & \textbf{Pregunta central} \\
\midrule
Clustering & ¿Hay grupos naturales? \\
Representación & ¿Puedo ver los datos en menos dims.? \\
Densidad/anomalías & ¿Hay puntos raros? \\
Autoencoders & ¿Puedo comprimir y reconstruir? \\
\bottomrule
\end{tabular}
}

\medskip
\textbf{Hilo conductor:} los cuatro bloques trabajan sin etiquetas $y_i$. La estructura emerge de los datos mismos. La elección del método depende de la pregunta que uno quiere responder.
\end{columns}
\end{frame}

\end{document}
